\section{Which open questions are worth pursuing next?}

The review supports a focused research program. More character tables or more equivalent one-sided criteria would add less value than proving one new arithmetic estimate. The existing analytic interfaces already accept pointwise, sparse, averaged, and selected filtered forms of control.

\subsection{Highest priority: a fixed nonabelian field}

Two complementary benchmarks are useful. The earlier conductor-23 example has the simplest well-developed analytic target: if $H$ is the splitting field of $T^3-T-1$, its Dedekind zeta function contains a degree-two Hecke factor over $\Q(\sqrt{-23})$ that is not a Dirichlet factor over $\Q$. The new real discriminant-229 example is particularly suited to extending the internal CM-field construction, because its Galois closure is totally real and can be adjoined without destroying the CM structure.\cite{priorMean,Cogdell}

For the conductor-23 field, write
\[
 -\frac{\zeta_H'}{\zeta_H}(s)=\sum_{n\ge1}\frac{a_H(n)}{n^s},
 \qquad
 \zeta_H(s)=\frac{\kappa_H}{s-1}+\kappa_H\gamma_H+O(s-1).
\]
The positive coefficients satisfy $0\le a_H(n)\le6\Lambda(n)$. Define
\[
 R_H(x)=\sum_{n\le x}a_H(n)\left(\frac1n-\frac1x\right)
         -\log x+1+\gamma_H.
\]
Its Mellin transform is
\[
 -\frac{\zeta_H'(s)/\zeta_H(s)}{s(s-1)}
 -\frac1{(s-1)^2}+\frac{1+\gamma_H}{s-1}.
\]
The singularity at $s=1$ is removable. A zero $\rho$ of multiplicity $m_\rho$ gives residue $m_\rho/[\rho(1-\rho)]$, which is positive at real zeros in $(0,1)$. Therefore
\begin{equation}\label{eq:target-mean}
 \int_X^{2X}\pos{R_H(x)}\,dx\ll_\varepsilon X^{1/2+\varepsilon}
\end{equation}
would exclude zeros to the right of $1/2$ and the functional equation would finish RH for this field. A square moment $\ll_\varepsilon X^\varepsilon$ also suffices. These are established reductions in the earlier reports, included here to specify the actual remaining arithmetic objective.\cite{priorMean}

The two filtered bounds of Corollary~\ref{cor:joint-half} are an alternative target if direct-sum theta constructions naturally provide the combinations at $x,2x,4x$. The positive-ramp theorem explains their pointwise content, but no conversion to a pointwise square-root estimate is required before using Landau.

For the discriminant-229 construction, the next useful theorem is an odd-cocycle bound over the nonabelian CM tower with the right uniformity in the tower level and weight, or a globally admissible replacement for the source's two coprime integral frames. Either would remove a real obstruction. A proof of mere finiteness, with uncontrolled dependence on the growing parameters, is insufficient for the final scaling argument.

\subsection{A nearby arithmetic classification: exponent seven}

The validated exponent-six report also supplies an exact upper bound
\[
 |D|\le4\cdot4600^7
 =174\,327\,062\,886\,400\,000\,000\,000\,000
\]
for a possible omitted exponent-seven field. Its rational explicit-formula checks were rerun as part of the same script. EKN's initial range contains 33 such fields; the residual interval remains unexcluded here.\cite{EKN,priorSix}

The exponent-six collision works because the forced number of split primes exceeds the number of possible coefficients $b$ in the norm equation. At exponent seven the same parameter choices do not give that inequality. A successful extension needs a sharper split-prime test, a stronger restriction on the coefficients, or a multi-equation elimination. Reusing the sixth-power certificate with an exponent changed in the code would not prove completeness.

This is an attractive finite arithmetic problem, because its target and certificate format are concrete. It uses the availability of quadratic GRH together with new integral arithmetic, rather than requiring a nonabelian RH extension first.

\subsection{Other candidates and their actual status}

\begin{longtable}{@{}p{.22\textwidth}p{.35\textwidth}p{.36\textwidth}@{}}
\toprule
Target & Relevant idea & Missing substantive result\\
\midrule
\endhead
Hecke GRH over a fixed nontrivial field & Prime divisibility encoded by embeddings; one-sided residue barrier & A matched height/index/rank inequality for the new coefficients.\\[4pt]
Nonabelian Dedekind RH & Complete-splitting selectors; positive regular-character coefficients & The arithmetic estimate, plus auxiliary-twist control where the chosen statistic needs it.\\[4pt]
Artin holomorphy & Integral periods, congruences, determinant ratios & Holomorphy cannot be inferred from a generic zero-free estimate or weakened internal attachment notion.\\[4pt]
Selmer or $K$-group vanishing & Local twist selection; denominator and Fitting-ideal bounds & A precise comparison with the desired fixed module, controlled coefficient ramification, and a separating certificate.\\[4pt]
Sparse equations for singular moduli & Size-versus-integrality; conjugate-root and determinant constraints & After the prior five-term completion, new six-term or uniform support-growth arguments; old residual lists are superseded.\\[4pt]
Exponent-seven class-group census & Positive explicit-formula tests and norm collisions & Exclusion of the residual interval left by the validated upper bound.\\
\bottomrule
\end{longtable}

The prior sparse-CM reports contain additional arguments, not merely a quotation of Riffaut's known conditional result. This session checked selected exact Newton-disk, reduced-form and confluent-determinant inputs, but did not replay all of their global polynomial-identification and residue-sieve certificates. Those classifications are therefore catalogued with that narrower validation status; the exponent-six census is the arithmetic completion fully revalidated here.

\subsection{A practical stopping criterion for each transfer}

A transfer should be called a solved new zero problem only after all of the following mathematical steps have been supplied: a statistic with correct poles; integral or geometric objects that genuinely encode that statistic; upper and lower bounds for the same normalized covolume; a rank and denominator estimate in the required range; and the analytic implication with no hidden cancellation. The present report completes some reusable local and analytic steps and one independent arithmetic classification. It leaves the new nonabelian covolume estimate open.
